\documentclass[11pt]{article}
\usepackage[margin=1in]{geometry}
\usepackage{lmodern}
\usepackage[T1]{fontenc}
\usepackage{amsmath,amssymb,amsthm}
\usepackage{booktabs}
\usepackage{hyperref}          % load hyperref last

% --- notation shortcuts ---
\newcommand{\R}{\mathbb{R}}
\newcommand{\x}{\mathbf{x}}
\newcommand{\w}{\mathbf{w}}
\DeclareMathOperator{\clamp}{clamp}
\DeclareMathOperator{\softplus}{softplus}
\DeclareMathOperator{\cumsum}{cumsum}
\DeclareMathOperator{\sg}{sg}          % stop-gradient
\newcommand{\ind}[1]{\mathbf{1}\!\left[#1\right]}
\newcommand{\radix}{\mathrm{radix}}

\title{Multi-Head LUT Detectors:\\ Anchor-Pairs, Hyperplane, and LIF}
\author{Anatoli Starostin}
\date{\today}

\begin{document}
\maketitle

\begin{abstract}
The Multi-Head LUT (MHL) family builds a differentiable function from a bank of small learned lookup tables
whose row addresses are produced by a bank of \emph{detectors} applied to the input. This note describes the
shared addressing--lookup--sum skeleton and three choices of detector: a fixed coordinate comparison
(anchor-pairs), a learned affine half-space (hyperplane), and a spiking leaky integrate-and-fire (LIF) neuron.
The first two discretize the input \emph{geometrically}; the third discretizes it \emph{temporally}, encoding
the input as spike latencies and reading out \emph{when} a neuron first fires. It is a mechanism/derivation
note only --- no experimental results are reported.
\end{abstract}

\section{The shared skeleton}
\label{sec:skeleton}

Let the input be $\x \in \R^{d_{\mathrm{in}}}$ (batch size $B$). A module has $H$ output heads and $P$ tables
per head, hence $T = H\,P$ lookup tables. Each table $t$ carries $D$ \emph{detectors}; each detector emits a
single integer \emph{digit}; the digits index a learned table of cells. The output is a tensor of shape
$(B, H, d_{\mathrm{out}})$.

Every MHL variant is the same four steps.

\paragraph{(1) Per-detector discretization.} Detector $d$ of table $t$ maps the input to a digit
\begin{equation}
  g_{t,d}(\x) \in \{0, 1, \dots, r_d - 1\},
\end{equation}
where $r_d$ is that detector's \emph{radix} (the number of symbols it may emit).

\paragraph{(2) Mixed-radix cell address.} The $D$ digits of a table are packed, row-major, into one cell index
\begin{equation}
  \mathrm{idx}_t \;=\; \sum_{d=0}^{D-1} g_{t,d}\,\radix_d,
  \qquad
  \radix_d \;=\; \prod_{d' > d} r_{d'} ,
  \label{eq:mixedradix}
\end{equation}
so the table has $\prod_{d} r_d$ rows (cells). When every $r_d = 2$ this is simply reading the digits as a
binary number, most-significant bit first.

\paragraph{(3) Table read.} Each table is a learned tensor $\mathrm{Table}_t \in \R^{\,\text{cells}\times
d_{\mathrm{out}}}$, and the addressed row is
\begin{equation}
  \mathrm{row}_t \;=\; \mathrm{Table}_t[\,\mathrm{idx}_t\,] \;\in\; \R^{d_{\mathrm{out}}}.
\end{equation}

\paragraph{(4) Head sum $\rightarrow$ output.} The $P$ tables of head $h$ are summed:
\begin{equation}
  \mathbf{y}_h \;=\; \sum_{p=1}^{P} \mathrm{row}_{(h,p)} \;\in\; \R^{d_{\mathrm{out}}},
  \qquad
  \mathbf{y} = (\mathbf{y}_1, \dots, \mathbf{y}_H) \in \R^{H \times d_{\mathrm{out}}}.
\end{equation}

The tables are the only memory; the detectors are the only thing that reads structure out of $\x$. The three
variants below differ \emph{only} in the choice of detector. Two of them (anchor-pairs, hyperplane) are
\emph{geometric}: a digit says which side of a boundary $\x$ lies on, with radix $r_d = 2$. The third (LIF) is
\emph{temporal}: a digit says \emph{when} a spiking neuron first fires, with radix $r_d = M$ time buckets.

\section{Anchor-pairs detectors (radix $2$, coordinate comparison)}
\label{sec:anchor}

The simplest detector is a fixed ordered pair of input coordinates $(a_d, b_d)$. Its digit is the sign of
their difference,
\begin{equation}
  g_{t,d}(\x) \;=\; \ind{\, x_{a_{t,d}} - x_{b_{t,d}} > \varepsilon \,} \;\in\; \{0,1\},
  \qquad \varepsilon \ge 0,
  \label{eq:anchor}
\end{equation}
where $\varepsilon$ is a comparison margin. Each table has $D$ such pairs; the $D$ bits form a binary index
into a table of $2^{D}$ rows via \eqref{eq:mixedradix}. The pairs are chosen once (balanced coverage over the
input coordinates) and held fixed; only the tables are learned in the base construction.

Because the boundary $x_a = x_b$ is a hard threshold, the digit is piecewise constant in $\x$ and the raw map
is non-differentiable. The module additionally exposes, per table, \emph{alternative} indices obtained by
flipping the bits with the smallest margins $|x_{a_d} - x_{b_d}|$, together with a per-bit \emph{uncertainty}
weight (for example inverse-$L_1$, $\propto 1/|x_{a_d} - x_{b_d}|$). These quantify which detector is closest
to flipping and are what a surrogate backward pass uses to route gradients to $\x$ and to neighbouring cells.
Geometrically each detector is the axis-pair hyperplane $x_a - x_b = 0$.

\section{Hyperplane detectors (radix $2$, learned affine)}
\label{sec:hyperplane}

This replaces the fixed coordinate pair by a learned affine half-space test. Detector $d$ of table $t$ owns a
weight row $\w_{t,d} \in \R^{d_{\mathrm{in}}}$ and a bias $b_{t,d} \in \R$, and emits
\begin{equation}
  g_{t,d}(\x) \;=\; \ind{\, \langle \w_{t,d}, \x \rangle + b_{t,d} > 0 \,} \;\in\; \{0,1\}.
  \label{eq:hyperplane}
\end{equation}
Again $D$ bits pack most-significant-first into a $2^{D}$-row table index, and the $P$ tables of a head are
summed. This strictly generalizes Section~\ref{sec:anchor}: initializing $\w_{t,d} = \mathbf{e}_{a} -
\mathbf{e}_{b}$ (a difference of standard basis vectors) and $b_{t,d} = 0$ gives $\langle \w_{t,d}, \x \rangle
= x_a - x_b$, exactly the anchor-pairs bit \eqref{eq:anchor}. The added freedom is that each detector's
boundary is now an arbitrary oriented hyperplane in input space rather than an axis-pair diagonal.

Training is end-to-end through the hard $\arg\max$/gather via a full-$K$ softmax surrogate: the forward value
is the hard table row, while the backward pass replaces each hard sign by a soft membership
$\sigma\!\big((\langle \w, \x\rangle + b)/\tau\big)$ and forms a soft distribution over the $2^{D}$ rows, so
gradients reach $\w$, $b$, the table weights, and $\x$. Two temperatures (a per-bit score temperature and a
row-select temperature) control the sharpness of that surrogate.

\section{LIF detectors (radix $M$, temporal / first-spike)}
\label{sec:lif}

Here a detector is a spiking leaky integrate-and-fire (LIF) neuron, and its digit is \emph{when} it first
fires, discretized into $M$ time buckets. A table carries $D$ such neurons, giving $M^{D}$ cells. The case
$D = 1$ is the ``bucket'' variant (one neuron, $M$ time buckets $\rightarrow M$ cells); $D > 1$ is the
``product'' variant (a mixed-radix address over several neurons $\rightarrow M^{D}$ cells).

\paragraph{(a) Latency coding.} The real input is first turned into input \emph{spike times}: each coordinate
$x_i$ fires once, earlier for larger $x_i$, through a fixed affine-then-clamp map onto a time window
$[0, t_w]$,
\begin{equation}
  \ell_i \;=\; \clamp\!\Big( t_w\big(\tfrac{1}{2} - \tfrac{3 x_i}{32}\big),\; 0,\; t_w \Big).
  \label{eq:latency}
\end{equation}
At $t_w = 32$ this is $\ell_i = \clamp(16 - 3 x_i,\, 0,\, 32)$: $x_i = 0$ fires at the window centre, the map
saturates at $x_i = \pm 16/3 \approx \pm 5.33$, and it is monotonically \emph{decreasing} (a larger input
fires earlier). The map has no learned parameters and expects roughly unit-variance input.

\paragraph{(b) Synapses, delays, arrivals.} Detector $d$ has one synapse per input coordinate, with a bounded
excitatory weight and a non-negative delay,
\begin{equation}
  w_{d,i} \;=\; w_{\max}\,\sigma\!\big(w^{\mathrm{raw}}_{d,i}\big) \in (0, w_{\max}),
  \qquad
  \delta_{d,i} \;=\; \clamp\!\big(\delta^{\mathrm{raw}}_{d,i},\, 0,\, t_w\big),
\end{equation}
where $\sigma$ is the logistic function. The arrival time of input $i$ at neuron $d$ is the input spike time
plus the synaptic delay,
\begin{equation}
  a_{d,i} \;=\; \ell_i + \delta_{d,i} \;\in\; [0, 2 t_w].
\end{equation}

\paragraph{(c) LIF membrane and first spike.} The neuron is a leaky integrator with time constant
$\tau_d = \softplus(\tau^{\mathrm{raw}}_d) + 1$. Its membrane potential is the sum of exponentially decaying
post-synaptic bumps from all arrivals up to time $t$,
\begin{equation}
  V_d(t) \;=\; \sum_{i:\, a_{d,i} \le t} w_{d,i}\,\exp\!\Big(-\tfrac{t - a_{d,i}}{\tau_d}\Big).
  \label{eq:membrane}
\end{equation}
Evaluated at the arrival instants (after sorting them in time), \eqref{eq:membrane} is computed in a single
pass by the factorization
\begin{equation}
  V \;=\; \exp\!\big(-a/\tau\big) \,\odot\, \cumsum\!\big( w \odot \exp(a/\tau) \big),
  \label{eq:cumsum}
\end{equation}
which is $O(N)$ in the number of synapses $N$. The neuron fires at its \emph{first-spike time}, the earliest
arrival at which the membrane crosses a fixed threshold $\theta = 1$,
\begin{equation}
  t^{\ast}_d \;=\; \min\big\{\, a_{d,i} : V_d(a_{d,i}) \ge \theta \,\big\},
  \qquad t^{\ast}_d = t_w \ \text{if it never crosses.}
  \label{eq:firstspike}
\end{equation}
A detector's output is thus a single scalar \emph{time} $t^{\ast}_d \in [0, t_w]$ --- the temporal analogue of
a hyperplane's scalar score $\langle \w, \x\rangle + b$.

\paragraph{(d) Temporal bucketing $\rightarrow$ digit.} Each detector owns $M-1$ trainable, strictly
increasing bucket boundaries, kept ordered by construction,
\begin{equation}
  \beta_{d,k} \;=\; \beta^{\mathrm{base}}_d + \sum_{j \le k} \softplus\!\big(\beta^{\mathrm{raw}}_{d,j}\big),
  \qquad \beta_{d,1} < \beta_{d,2} < \dots < \beta_{d,M-1}.
\end{equation}
The digit counts how many boundaries the first-spike time has passed,
\begin{equation}
  g_{t,d}(\x) \;=\; \#\{\, k : t^{\ast}_d \ge \beta_{d,k} \,\}
             \;=\; \sum_{k=1}^{M-1} \ind{\, t^{\ast}_d \ge \beta_{d,k} \,} \;\in\; \{0,\dots,M-1\}.
  \label{eq:lifdigit}
\end{equation}
This is a radix-$M$ digit (versus the radix-$2$ bit of Sections~\ref{sec:anchor}--\ref{sec:hyperplane}). Steps
(a)--(d) produce one digit per detector; the digits then run through the shared mixed-radix address, table
read, and head sum of Section~\ref{sec:skeleton}.

\paragraph{(e) Differentiable training.} The hard path --- the $\arg\min$ first-spike \eqref{eq:firstspike} and
the boundary count \eqref{eq:lifdigit} --- is piecewise constant. Training uses a straight-through estimator:
the forward value is the hard cell read $\mathbf{y}_{\mathrm{hard}}$ (which also carries the gradient to the
selected table row), while a parallel soft readout $\mathbf{y}_{\mathrm{addr}}$ supplies the address gradient,
combined so the forward stays exactly hard,
\begin{equation}
  \mathbf{y} \;=\; \mathbf{y}_{\mathrm{hard}} + \mathbf{y}_{\mathrm{addr}} - \sg\!\big(\mathbf{y}_{\mathrm{addr}}\big),
  \label{eq:st}
\end{equation}
where $\sg(\cdot)$ is the stop-gradient operator. The soft readout replaces the two hard steps by smooth
surrogates: a soft first-spike time from per-arrival crossing (``escape'') probabilities
$c_k = \sigma\!\big((V_k - \theta)/T_{\mathrm{cross}}\big)$ with survival $\prod_{j<k}(1 - c_j)$, giving an
expected spike time; and a soft bucket partition over the $M$ buckets from differences of
$\sigma\!\big((t - \beta_{d,k})/T_{\mathrm{bkt}}\big)$. The soft distribution over cells is read against a
\emph{detached} copy of the table, so that \eqref{eq:st} sends value and table-weight gradients through the
hard read while the soft read sends address gradients to the weights, delays, $\tau$, and boundaries. Two
temperatures $T_{\mathrm{cross}}, T_{\mathrm{bkt}}$ set the softness (they may be frozen to $1$). Inference
uses the pure hard path with no soft computation.

\section{How the three connect, and where LIF differs}
\label{sec:connect}

Side by side, the variants are one backbone with three detector choices (Table~\ref{tab:variants}); everything
from the digit onward --- the mixed-radix index \eqref{eq:mixedradix}, the table read, and the head sum --- is
identical.

\begin{table}[h]
  \centering
  \begin{tabular}{@{}llll@{}}
    \toprule
    Variant & Detector & Digit $g_d$ & Radix $r_d$ \\
    \midrule
    Anchor-pairs & fixed coordinate pair        & $\ind{x_a - x_b > 0}$             & $2$ \\
    Hyperplane   & learned affine half-space    & $\ind{\langle\w,\x\rangle + b > 0}$ & $2$ \\
    LIF          & spiking LIF neuron, 1st spike & $\#\{k : t^{\ast} \ge \beta_k\}$  & $M$ \\
    \bottomrule
  \end{tabular}
  \caption{The three detectors share the addressing and table-read backbone and differ only in how a digit is
  produced. A table then has $r_d^{\,D}$ cells: $2^{D}$ for the geometric variants, $M^{D}$ for LIF.}
  \label{tab:variants}
\end{table}

The geometric variants discretize the input \emph{spatially}: a detector answers ``which side of a boundary is
$\x$ on?'', the boundary being a coordinate-pair diagonal (anchor-pairs) or a general learned hyperplane
(hyperplane, which contains anchor-pairs as an initialization). Each detector yields one bit, and a table read
is a lookup on a sign pattern.

The LIF variant discretizes the input \emph{temporally}: $\x$ becomes a pattern of spike latencies
\eqref{eq:latency}, a LIF neuron integrates delayed, weighted, leaky charge \eqref{eq:membrane}, and the
detector reports \emph{when} it first fires \eqref{eq:firstspike}, sorted into $M$ learned time buckets
\eqref{eq:lifdigit}. Two consequences distinguish it from the geometric detectors.
\begin{itemize}
  \item \textbf{Richer per-detector code.} A geometric detector emits one bit; a LIF detector emits one of $M$
    symbols (a radix-$M$ digit), so a single neuron already splits the input into $M$ regions rather than $2$.
  \item \textbf{Dynamics instead of a static boundary.} The ``boundary'' is no longer a fixed hyperplane but
    the neuron's temporal dynamics --- per-synapse delays, a leak time constant $\tau$, and a threshold
    $\theta$ --- which decide both \emph{whether} and \emph{when} it fires. In the degenerate limit of no
    leak, no delays, and a single readout instant, threshold crossing on the summed weighted latencies reduces
    to a linear threshold on the $\ell_i$, i.e.\ an affine test on $\x$. The LIF detector therefore
    \emph{contains} the hyperplane detector as a special case and adds temporal structure (delays, leak,
    multiple time buckets) that a single hyperplane cannot express.
\end{itemize}

In short: the same lookup-table backbone throughout, but the detector moves from \emph{geometry} (which side
of a plane) to \emph{time} (when a spiking neuron fires), trading a binary spatial test for a multi-valued
temporal one.

\end{document}
